\section{Related Work}
\label{sec:related}

% -----------------------------------------------------------------------------
\subsection{Medical Image and Polyp Segmentation}
Automated lesion segmentation from colonoscopy video frames plays a vital
role in computer-aided detection (CADe) and diagnosis (CADx)
systems~\cite{siegel2023colorectal, leufkens2012factors}. Beginning with
the foundational U-Net architecture~\cite{ronneberger2015unet} and its
multi-scale variants such as UNet++~\cite{zhou2018unetplusplus} and
ResUNet~\cite{zhang2018road}, deep convolutional networks have established
standard benchmarks in biomedical segmentation. To better capture subtle
boundary transitions, attention-guided models have been proposed, including
Attention U-Net~\cite{oktay2018attention}, PraNet~\cite{fan2020pranet}
(which utilizes parallel reverse attention to model polyp contours), and
HarDNet-MSEG~\cite{huang2021hardnet}. More recently, vision transformers
such as TransUNet~\cite{chen2021transunet}, Swin-Unet~\cite{cao2022swin},
and Polyp-PVT~\cite{dong2021polyp} have leveraged self-attention to model
long-range spatial dependencies.

However, standard feedforward architectures remain inherently single-pass:
intermediate representations cannot be iteratively re-examined or refined
in light of downstream contextual decisions. In sessile polyp segmentation,
where lesions are flat and visual margins are ambiguous, single-pass
feedforward models often struggle to balance sensitivity against excessive
background false alarms.

% -----------------------------------------------------------------------------
\subsection{Feedback and Recurrent Refinement Networks}
In biological vision, anatomical feedback projections from higher cortical
areas to early visual processing stages outnumber feedforward projections
by a significant margin~\cite{felleman1991distributed, gilbert2013top}.
Inspired by this principle, recurrent feedback networks have been developed
to simulate iterative cognitive refinement in artificial neural
networks~\cite{liang2015recurrent, zamir2017feedback}. In general computer
vision, Feedback Networks~\cite{zamir2017feedback} and recurrent
convolutional networks (RCNN)~\cite{pinheiro2014recurrent} demonstrate that
re-routing higher-level predictions to lower-level feature extractors
improves object localization and contextual disambiguation.

In medical imaging, Recurrent U-Net~\cite{alom2018recurrent} and
particularly the Feature Attention Network (FANet)~\cite{tomar2021fanet}
introduced cross-epoch feedback loops. FANet introduces the \texttt{MixPool}
module, reinjecting the previous epoch's binary prediction mask directly
into the encoder and decoder stages to guide feature extraction.

\textbf{The Overlooked Research Gap:}
Existing literature in recurrent segmentation operates on the heuristic
assumption that reinjecting prediction masks into intermediate layers
automatically encourages iterative refinement. Crucially, these studies
treat the recurrent coupling as a black box, completely overlooking the
gradient dynamics of the feedback interface. In this work, we demonstrate
that hard-thresholded mask injection creates a catastrophic
\textit{Feedback Trap}: it attenuates gradient backpropagation by over
$79\%$, paralyzes internal attention learning, and causes early encoder
representations to drift uncontrollably.

% -----------------------------------------------------------------------------
\subsection{Loss Formulations for Imbalanced Medical Segmentation}
Segmentation of small, irregular lesions is heavily impacted by class
imbalance, where background pixels vastly outnumber foreground lesion
pixels~\cite{salehi2017tversky}. Standard Cross-Entropy often biases
predictions toward the background, prompting the widespread adoption of
the Dice loss~\cite{milletari2016vnet} and soft Jaccard
loss~\cite{rahman2016optimizing}.

To explicitly penalize false-positive over-segmentation or false-negative
misses, asymmetric and boundary-aware loss formulations have emerged. The
Tversky loss~\cite{salehi2017tversky} generalizes the Dice index by
introducing weighting parameters $\alpha$ and $\beta$ to independently
penalize false positives and false negatives. The Focal Tversky
loss~\cite{abraham2019novel} further modulates hard examples, while
Asymmetric Unified Focal loss~\cite{yeung2022unified} adaptively balances
precision and recall.

\textbf{Our Architectural Perspective:}
Contemporary research overwhelmingly attempts to suppress false positives
by designing increasingly intricate loss functions. However, our factorial
experiments reveal that such loss engineering is entirely neutralized when
paired with conventional recurrent feedback loops. Rather than proposing
another loss formulation, our work addresses the structural root of the
problem: by redesigning the recurrent gating mechanism to eliminate zero
gradients and detach corruptive feedback paths, we \textbf{unleash the
latent corrective power of existing asymmetric loss functions}, enabling
them to suppress over-segmentation effectively.
